\section{Proof Sketch}\label{sec:proof-sketch}

The coding scheme of \citet[Theorem~16]{lau2026orderoptimal} first produces
an interval of length at most \(100\sigma\); taking its midpoint gives an
always-defined center \(c\) with \(|c-\mu|\leq50\sigma\) except on an event
of probability at most \(\delta/4\).  Recentring the \(k\)-moment at this
actual decoder output costs only a \(k\)-dependent constant
(Lemma~\ref{lem:step-002-recentered-moment}).

At each dyadic scale, one of four shifts puts \(c\) a distance at least
\(3h_j/8\) from the selected cell boundary
(Proposition~\ref{prop:step-004-cell-margin}).  The encoder queries every
randomly selected scale-offset entry without knowing \(c\); the decoder
retains the matching entry afterward.  Uniform dithering turns its centered
bit into an exact selected-digit difference
(Proposition~\ref{prop:step-005-dither-identities}), and these differences
telescope to \(X-c\) up to a bottom quantization residual and a top
truncation residual (Proposition~\ref{prop:step-006-residual-interface}).

The cell margin makes a digit identically zero until
\(|X-c|>3h_j/8\).  Summing this activation constraint pathwise before taking
expectations gives a count-free fine-scale bound and a \(k\)-moment-compatible
coarse bound (Propositions~\ref{prop:step-007-fine-ledger} and
\ref{prop:step-007-coarse-ledger}).  Consequently the conditional variance
has the three regimes \(\sigma^2\), \(\sigma^2\log(\sigma/\epsilon)\), and
\(\sigma^kH^{2-k}\), with only one logarithm at \(k=2\)
(Proposition~\ref{prop:step-010-three-regime}).

The two residual expectations total at most \(\epsilon/4\)
(Proposition~\ref{prop:step-011-bias-certificate}).  A fixed
median-of-means allocation estimates the exact telescope mean to accuracy
\(\epsilon/2\), after which the localization and refinement failure
probabilities are integrated under their common joint law
(Proposition~\ref{prop:step-013-unconditional-pac}).  Finally, the exact
ceilings, confidence factors, and localization cost are absorbed into the
three-regime rate in the Rate Specialization Bridge,
Proposition~\ref{prop:step-014-rate-bridge}.
